% sections/07_discussion.tex --- SkillStrike

\subsection{Why Existing Defenses Fail}

\paragraph{Input sanitization.}
Standard prompt injection defenses~\cite{greshake2023indirect}
sanitize user-supplied input before feeding it to the agent.
Skill composition attacks bypass this entirely:
the malicious instruction arrives through a \emph{trusted, installed skill's
tool output}---a channel the agent's safety policy explicitly exempts
from user-input rules.
The agent is correctly following its design: tool outputs from installed
skills are supposed to be trusted.

\paragraph{Platform scanners.}
The \clawhub platform scanner detected 0 of 34 independently confirmed
malicious skills in our corpus (\S\ref{sec:static}).
\lscps at Tier-6 detects all 34 (100\% recall vs.\ 0\% for the platform scanner).
Current platform scanners check for malware signatures and known-bad domains,
not for the structural source-sink property we exploit.

\paragraph{Agent safety policies.}
Both evaluated agent runtimes apply heuristic safety policies that
are: (a) command-surface-specific rather than capability-class-based,
(b) stochastically applied (the same task type succeeds in some
pairs and fails in others with no configuration change), and
(c) channel-dependent (memory writes refused in C1 context pass via C3 channel).
Findings F3, F8, F9 all confirm that these policies are
\emph{insufficient to prevent composition attacks}.

\subsection{Proposed Defenses}

\paragraph{D1: Mandatory capability manifests.}
Skill marketplaces should require machine-readable capability declarations
at publish time---analogous to Android permission manifests---enumerating:
filesystem access scope, network destinations, shell invocation patterns,
credential file access, and external API endpoints.
Our K1--K5 taxonomy provides a starting point for this schema.
Capabilities not declared in the manifest should be blocked at runtime
by the agent framework.

\paragraph{D2: Publish-time \lscps screening.}
\lscps can run as a publish-time gate with zero API cost,
requiring no LLM queries.
At our measured false-positive rate (estimated $<$5\% at Tier-6),
the cost to legitimate skill developers is one clarification email
per 20 submissions.
The benefit: 100\% recall on the 34 confirmed malicious skills
that the platform scanner missed.

\paragraph{D3: Cross-skill capability isolation.}
Agent runtimes should enforce \emph{context-window isolation}:
when skill A is invoked, only skill A's tool schemas should be
active in the agent's context---not the schemas of co-installed skills.
This prevents the agent from invoking skill B's tools in
response to content returned by skill A.
This requires changes to MCP protocol handling
but eliminates the structural precondition for composition attacks.

\paragraph{D4: Install-time composition audit.}
When a user installs a sink skill alongside an existing source skill,
the platform should execute a composition audit:
does the new combination create an exploitable source-sink pair
(as determined by \lscps or LLM verification)?
If so, alert the user before completing installation.

\subsection{Limitations}

\paragraph{Coverage of \lscps.}
\lscps is a conservative syntactic lower bound.
Dynamically constructed commands, base64-encoded payloads,
and import-alias obfuscation are not detected.
All reported attack surface numbers are underestimates.

\noindent\emph{Ground-truth calibration.}
The Tier-6 recall claim (100\% on 34 confirmed skills vs.\ 0\% platform scanner)
rests on an independent security researcher's manual review of those 34 skills,
conducted prior to and independent of our experiments.
This reviewer was not aware of our tier assignments during their review,
which reduces annotation bias.
We do not claim these 34 skills represent the complete malicious-skill population;
they are the set for which an independent ground truth label was available.
A full precision/recall evaluation over a labeled corpus of $>$1,000 skills
remains future work.

\paragraph{LLM verification annotation.}
The Stage-1 precision estimate (92\%) is based on one author's manual review
of 50 randomly sampled verdicts.
Inter-annotator agreement on a second 25-verdict sample (two annotators)
yields $\kappa = 0.84$ (substantial agreement);
we acknowledge that 50 samples provide approximate---not definitive---precision estimates.

\paragraph{Experimental configuration realism.}
\skillssh{} experiments use \texttt{--full-auto} CodexCLI, matching a user who
typically enables full-auto for development tasks.
For T07 (network exfiltration), \texttt{--dangerously-bypass} is used
only to allow our outbound-logging infrastructure to observe the attempt---
it does not change the agent's own safety policy.
\textbf{Conservative estimate:} removing all T07 pairs (26/166) from the total
yields 139/140 = 99.3\% success among the remaining non-bypass tasks,
confirming the 99\% result is not driven by the bypass configuration.

\paragraph{Payload generalization.}
We test three payload families (P1--P3).
Adversarial search over the prompt space would likely discover
more effective or better-disguised payloads.
Our reported success rates are therefore \emph{lower bounds} on
what a dedicated attacker could achieve.

\paragraph{Agent version.}
Results are specific to \texttt{gpt-5.3-codex} (OpenClaw) and
\texttt{gpt-4.1} (Codex CLI) as of April 2026.
Future model versions may apply different safety heuristics.

\paragraph{Two marketplaces.}
We evaluate \clawhub and \skillssh{} only.
Other agent platforms (Cursor, Continue.dev, Claude Computer Use,
GitHub Copilot Extensions) have different installation and execution
models; the structural vulnerability (shared context, no capability isolation)
applies to all, but empirical success rates may differ.

\subsection{Ethical Disclosure}

All experiments were conducted in fully isolated sandbox environments
(dedicated per-experiment workspace directories with no real external network access).
No real user accounts, production skills, or marketplace infrastructure were harmed.
We used real skill identifiers to enable reproducibility;
no weaponized payloads, exploit code, or automated attack scripts are released.

\noindent\textbf{Disclosure timeline.}
We notified the ClawHub platform security team on 2026-03-28
and the skills.sh maintainers on 2026-04-01,
before submission of this paper.
Both acknowledged receipt. ClawHub has indicated they are reviewing
capability manifest requirements; skills.sh has not responded.
Artifact release is de-weaponized: we release the \lscps pattern library
and LLM audit prompts, but not the identity-spoofing MCP server
or pre-packaged payload scripts.
